% !PRINTED none
% Chapter-opening page: no running head/page number is printed at the top of this
% page; instead the page carries the Abschnitt title block, then the §.115 heading
% and the start of the section text. (Weber suppresses the folio number on
% Abschnitt-opening pages.)

\begin{center}
\gesperrt{Vierzehnter Abschnitt.}

\gesperrt{Gruppen linearer ternärer Substitutionen.}

\rule{0.3\linewidth}{0.4pt}

\S.\,115.

\gesperrt{Ternäre lineare Substitutionsgruppe vom 168\textsuperscript{sten} Grade.}
\end{center}

Die Frage nach der Gesammtheit aller möglichen endlichen
Gruppen linearer Substitutionen von mehreren, insbesondere
derer von drei und vier Dimensionen, ist von \gesperrt{C. Jordan} behandelt,
der, ähnlich wie es für die binären Substitutionen geschehen,
eine endliche Anzahl von Typen solcher Gruppen aufgestellt hat\footnote{\gesperrt{C. Jordan}, Mémoire sur les équations différentielles linéaires à
intégrale algébrique. Crelle's Journal für Mathematik, Bd. 84 (1878). Sur
la détermination des groupes d'ordre fini etc. Atti della R. Accademia di
Napoli, Vol. VIII (1879). \gesperrt{Valentiner}, Kj\"ob. Skrift (6) V (1889).}. Diese allgemeinen Untersuchungen können wir
hier nicht verfolgen, wir begnügen uns, ein Beispiel einer ternären
Gruppe ausführlicher zu betrachten.

In dem Abschnitte über die Congruenzgruppen (\S.\,66) haben
wir eine Reihe \gesperrt{einfacher} Gruppen kennen gelernt, von denen
sich die erste, vom Grade 60, als isomorph mit der Ikosaëdergruppe
erwiesen hat. Wir wollen nun die nächste dieser einfachen
Gruppen vom Grade 168 betrachten, die wir mit $L_7$ bezeichnet
haben, und versuchen, eine mit dieser isomorphe, ternäre, lineare
Substitutionsgruppe zu construiren.

Dazu führt uns das Theorem \S.\,72, I., wenn wir vier Elemente
$\tau, \chi, \omega, \Theta$ aufsuchen, die bei der Composition den Bedingungen
dieses Theorems genügen.

Wir wollen zunächst eine Substitution $\tau$ in der Normalform
annehmen
\weq{1}{\tau = \begin{pmatrix} \varepsilon_1, & 0, & 0 \\ 0, & \varepsilon_2, & 0 \\ 0, & 0, & \varepsilon_3 \end{pmatrix},}
%CONT>>
